\section{Conformance and Evaluation}

WorldEpisode should be evaluated as an executable standard, not only as a schema.

\paragraph{RQ1: semantic preservation across bindings.}
Convert the same dataset through LeRobot, WorldEpisode, Rerun, NCore, MCAP, OpenUSD, and glTF
Gaussian assets. Measure field retention, semantic retention, deterministic round-trip equality,
externalized fields, approximated fields, and declared loss.

\paragraph{RQ2: physical-coherence fault detection.}
Inject controlled failures: wrong transform direction, millimeters interpreted as meters,
quaternion-order mismatch, timestamp offset, stale calibration, action frame mislabeled, collision
mesh scale mismatch, reused entity id, or world revision changed without a new hash. Measure
validator precision, recall, and diagnosis quality.

\paragraph{RQ3: cross-simulator replay.}
Replay the same demonstrations in at least two simulators. Measure end-effector RMSE, object
trajectory RMSE, grasp-state agreement, contact-event precision/recall/F1, final-state pose error,
task-outcome agreement, and success-rate rank correlation across policy checkpoints.

\paragraph{RQ4: downstream robustness.}
Train with observations/actions alone, with unstructured 3D side files, and with full WorldEpisode
identity/world alignment. Evaluate shifted camera extrinsics, unseen viewpoints, object displacement,
background replacement, distractors, occlusion, lighting changes, novel object instances, and held-out
worlds.

\paragraph{RQ5: lineage-safe splits.}
Compare random episode splits against task-, world-, entity-, and reconstruction-lineage-disjoint
splits. A substantial performance drop under lineage-safe splits would show that current dataset
abstractions overestimate generalization.

The minimum credible reference release should include a small golden conformance corpus: valid
minimal packages, valid advanced packages, intentionally invalid packages, lossy-conversion examples,
and cross-simulator replay fixtures.
